%%%%%%%%%%  scan idx 308  |  folio 295  %%%%%%%%%%
On en
déduit en particulier (en faisant \(X=Y\), \(f=\Id\)) que \(K_{\bullet}(g)\) est un \(K^{\bullet}(Y)\)-module. D'autre part, si \(f\) est parfait, le produit par \(\cl(\mathcal O_X)\in K^{\bullet}(f)\) définit des homomorphismes \(K^{\bullet}(Y)\)-linéaires
\[
 f^*:K_{\bullet}(Y/S)\longrightarrow K_{\bullet}(X/S),
\]
et
\[
 f^*:K^{\bullet}(Y/S)\longrightarrow K^{\bullet}(X/S).
\]

\item Supposons \(S\) quasi-compact, \(X\) et \(Y\) de type fini sur \(S\), et \(f\) propre. Alors, sous réserve de validité de la conjecture de finitude (III 2.1), donc en particulier si \(Y\) est noethérien ou si \(f\) est projectif, le foncteur \(Rf_*\) envoie \(D^b(h)_{\coh}\) dans \(D^b(g)_{\coh}\), donc définit un homomorphisme
\[
 f_*:K_{\bullet}(X/S)\longrightarrow K_{\bullet}(Y/S).
\]
Si en outre \(f\) est parfait, \(Rf_*\) envoie \(D(h)_{\parf}\) dans \(D(g)_{\parf}\) (III 4.8), et définit un homomorphisme
\[
 f_*:K^{\bullet}(X/S)\longrightarrow K^{\bullet}(Y/S).
\]
De la formule de projection (III 3.7) on déduit la formule
\[
 f_*(a x)=a f_*(x)
\]
pour \(f\) propre, \(a\in K^{\bullet}(S)\), \(x\in K_{\bullet}(X)\) (resp. \(f\) propre et parfait, \(a\in K_{\bullet}(S)\), \(x\in K^{\bullet}(X/S)\)).

\item Supposons \(S\) noethérien, \(X\) et \(Y\) de type fini sur \(S\), et \(f\) compactifiable et parfait. Alors (III 4.9.1) le foncteur \(f^!\) envoie \(D(g)_{\parf}\) dans \(D(h)_{\parf}\), et définit un homomorphisme
\[
 f^!:K^{\bullet}(Y/S)\longrightarrow K^{\bullet}(X/S).
\]
\end{enumerate}

\textbf{3.3.3.} Sous les hypothèses de (III 4.7), le foncteur \(\otimes_S^{\mathrm L}\) définit des accouplements
\[
 D(f_1)_{\parf}\times D^b(f_2)_{\coh}\longrightarrow D^b(f)_{\coh},
\]
\[
 D(f_1)_{\parf}\times D(f_2)_{\parf}\longrightarrow D(f)_{\parf},
\]
qui induisent des homomorphismes
\[
 K^{\bullet}(f_1)\otimes K_{\bullet}(f_2)\longrightarrow K_{\bullet}(f),
\]
\[
 K^{\bullet}(f_1)\otimes K^{\bullet}(f_2)\longrightarrow K^{\bullet}(f).
\]
